% Sección a cargo de: Gilces Carranza José Ignacio
\section{Declaración de uso de inteligencia artificial}\label{sec:ia}

El equipo declara el uso de una herramienta de inteligencia artificial durante la elaboración de esta entrega, según la norma de entrega de la asignatura. La Tabla~\ref{tab:ia} resume la herramienta, el propósito de cada uso y la forma en que se verificó el resultado.

\begin{table}[H]
\centering
\caption{Uso de inteligencia artificial en la Entrega 1.}
\label{tab:ia}
\small
\begin{tabularx}{\textwidth}{@{}L{3.0cm}X@{}}
\toprule
\textbf{Campo} & \textbf{Detalle} \\
\midrule
Herramienta & Claude Code, aplicación de escritorio de Anthropic. \\
Versión del modelo & Claude Opus 5.5 (\texttt{claude-opus-5-5}). \\
Periodo de uso & 19 al 21 de noviembre de 2026. \\
Insumos entregados por el equipo & Guía y rúbrica TA-PFC-E1, informe de la práctica de la Unidad~I, guía de entrevistas ACTA-ELI-01, transcripciones de las tres entrevistas del 19 de noviembre y formulario de consentimiento PFC-ISR401-POSIA-A03. \\
\bottomrule
\end{tabularx}
\end{table}

\subsection*{Propósito de cada uso}

\begin{enumerate}[label=\arabic*., leftmargin=*]
\item \textbf{Planificación.} Reparto de la actividad entre los integrantes y cronograma de commits.
\item \textbf{Estructura del repositorio.} Organización de carpetas, plantilla de \texttt{main.tex}, \texttt{README.md}, \texttt{.gitignore} y formato del manifiesto de evidencias.
\item \textbf{Paso a \LaTeX{} de material propio.} Conversión de la guía ACTA-ELI-01 y del formulario de consentimiento del equipo a los anexos del documento.
\item \textbf{Borradores de redacción.} Primeras versiones de las secciones de visión, alcance, partes interesadas, contexto, elicitación, ética, gestión, riesgos y conclusiones, elaboradas a partir de las transcripciones y de los documentos del equipo.
\item \textbf{Modelos.} Código TikZ de la matriz poder--interés y del diagrama de contexto, con la ubicación de los elementos definida a partir de las tablas del documento.
\item \textbf{Revisión de consistencia.} Detección de diferencias entre el texto y las evidencias, como la custodia descrita frente a la del formulario firmado y el correo pedido sin autorización en una entrevista.
\end{enumerate}

\subsection*{Lo que no se hizo con inteligencia artificial}

Las entrevistas, sus grabaciones y transcripciones, el formulario de consentimiento, la guía ACTA-ELI-01, la carta de aceptación y el contacto con la organización son trabajo de campo del equipo. Ningún hallazgo del documento proviene de la herramienta: todos se toman de las transcripciones de las sesiones ACTA-ELI-CAJ-01, ACTA-ELI-CAJ-02 y ACTA-ELI-ADM-01.

\subsection*{Verificación realizada por el equipo}

\begin{enumerate}[label=\arabic*., leftmargin=*]
\item Cada integrante revisó y ajustó el borrador de su sección y comprobó que cada hallazgo citado aparece en la transcripción indicada.
\item Se comprobaron en \url{https://doi.org} los DOI de las referencias citadas y se compararon autores, título, año y páginas con la obra real.
\item Se comparó el anexo de consentimiento con el formulario firmado por los participantes.
\item Se compiló el documento siguiendo el \texttt{README.md} y se revisó que no hubiera referencias sin resolver.
\end{enumerate}

Los integrantes asumen la responsabilidad del contenido entregado y pueden explicarlo y modificarlo cuando se les solicite.
